% ===================== III. Sivers shift from amplitudes =====================
\section{Sivers shift}

% Slide 10
\begin{frame}{From the correlator to invariant amplitudes}
  \begin{itemize}
    \item Decompose the $\bvec$-space correlator into Lorentz-invariant amplitudes. For $\Gamma=\gamma^+$, two of them carry $f_1$ and the Sivers function:
  \end{itemize}
  \[ \frac{1}{2 P^+}\widetilde \Phi^{[\gamma^+]}_{\unsub} = \textcolor{red}{\tAmp_{2B}} + i \mN \epsilon_{ij}\, \vect{\bvec}_i \vect{S}_j\, \textcolor{red}{\tAmp_{12B}}, \qquad
     \tAmp_{2B} \equiv \tAmp_2 + R(\zetahat^2) \tBmp_1, \quad \tAmp_{12B} \equiv \tAmp_{12} - R(\zetahat^2) \tBmp_8 \]
  \begin{itemize}
    \item The amplitudes depend only on invariants, so they can be computed in a convenient lattice frame:
  \end{itemize}
  \[ \tAmp_i\big(\bvec^2,\ \bvec \tcdot P,\ (\bvec \tcdot P) R(\zetahat^2)/\mN^2,\ -1/(\mN\zetahat)^2,\ \eta v \tcdot P\big) \]
  \begin{itemize}
    \item $\zetahat$: Collins--Soper parameter (how close $v$ is to the light cone). \quad $\eta v\tcdot P\to\pm\infty$: SIDIS / Drell--Yan staple.
  \end{itemize}
  \keymsg{At $\bvec\tcdot P=0$: \quad $\displaystyle\langle k_y\rangle_{TU}(\vprp{b}^2) = -\mN\,\frac{\tAmp_{12B}(-\vprp{b}^2,0,0,\ldots)}{\tAmp_{2B}(-\vprp{b}^2,0,0,\ldots)}$ \quad --- the $x$-\textbf{integrated} Sivers shift}
\end{frame}

% Slide 11
\begin{frame}{Key idea 1: $x$ is hiding in $\bvec\tcdot P$}
  \begin{itemize}
    \item $x$ and $\bvec\tcdot P$ are Fourier conjugates. Keep $\bvec\tcdot P$ as a variable and transform numerator and denominator:
  \end{itemize}
  \begin{block}{Generalized Sivers shift as a function of $x$}
  \[ \langle \vect{k}_y \rangle_{TU}(\vprp{\bvec}^2,\textcolor{blue}{x},\zetahat,\eta v \tcdot P)
     \equiv \mN \frac{\tilde f_{1T}^{\perp(1)}}{\tilde f_1^{(0)}}
     = -\mN \frac{\int d(\bvec\cdot P)\,e^{ix(\bvec\cdot P)}\, \textcolor{red}{\tAmp_{12B}}(\bvec^2,\bvec \tcdot P,\ldots)}{\int d(\bvec\cdot P)\,e^{ix(\bvec\cdot P)}\,\textcolor{red}{\tAmp_{2B}}(\bvec^2,\bvec \tcdot P,\ldots)} \]
  \end{block}
  \begin{itemize}
    \item Taylor coefficients in $\bvec\tcdot P$ are the $x$-moments at fixed $|\vprp b|$; the value at $\bvec\tcdot P=0$ is $\int dx$.
    \item Under $\Gamma$-independent multiplicative renormalization, the soft factor and $Z$ cancel in the amplitude ratio at fixed geometry $(\bvec, P, \eta v)$. We keep this ratio structure and scan $\bvec\tcdot P$.
  \end{itemize}
\end{frame}

% Slide 12
\begin{frame}{What the Fourier transform needs}
  \begin{itemize}
    \item Hermiticity: $\tAmp^{\text{Re}}$ is even, $\tAmp^{\text{Im}}$ is odd in $\bvec\tcdot P$. They separate quark and antiquark combinations:
  \end{itemize}
  \[ \myRe\,\tAmp_{2B} = \tfrac{\tilde{\softf}}{2}\!\int_0^1\! dx\, \cos(x\,\bvec\tcdot P)\big[\tilde f_1^q - \tilde f_1^{\bar q}\big], \qquad
     \myIm\,\tAmp_{2B} = -\tfrac{\tilde{\softf}}{2}\!\int_0^1\! dx\, \sin(x\,\bvec\tcdot P)\big[\tilde f_1^q + \tilde f_1^{\bar q}\big] \]
  \[ \Rightarrow\quad \langle \vect{k}_y \rangle_{TU}(x) = -\mN \frac{\tAmp_{12B}^{[x]\text{Re}}-\tAmp_{12B}^{[x]\text{Im}}}{\tAmp_{2B}^{[x]\text{Re}}-\tAmp_{2B}^{[x]\text{Im}}},
     \qquad \begin{aligned} \tAmp^{[x]\text{Re}} &= \!\textstyle\int\! d(\bvec\tcdot P)\cos(x\,\bvec\tcdot P)\,\tAmp^{\text{Re}} \\ \tAmp^{[x]\text{Im}} &= \!\textstyle\int\! d(\bvec\tcdot P)\sin(x\,\bvec\tcdot P)\,\tAmp^{\text{Im}} \end{aligned} \]
  \begin{itemize}
    \item $\tAmp_{12B}^{\text{Im}}$ is consistent with zero in our data and is dropped $\Rightarrow$ three amplitudes: $\tAmp_{2B}^{\text{Re}},\ \tAmp_{2B}^{\text{Im}},\ \tAmp_{12B}^{\text{Re}}$.
  \end{itemize}
  \begin{alertblock}{Two requirements to remember}
    \begin{enumerate}
      \item The amplitude must \textbf{decay in $\bvec\tcdot P$} fast enough for the integral to exist.
      \item The lattice covers only a \textbf{finite window} of $\bvec\tcdot P$; the transform needs all of it, so a fitted model supplies the tail. Small $x$ is the least constrained region.
    \end{enumerate}
  \end{alertblock}
\end{frame}
